\documentclass[11pt]{article}
\usepackage[margin=1in]{geometry}
\usepackage{amsmath,amssymb}
\usepackage{booktabs}
\usepackage{hyperref}
\usepackage{listings}
\usepackage{xcolor}
\lstset{basicstyle=\ttfamily\footnotesize, breaklines=true, columns=fullflexible}
\title{Independent Replication Report:\\ \textit{A modular functor which is universal for quantum computation}\\ (Freedman, Larsen, Wang, arXiv:quant-ph/0001108)}
\author{QC-200 Replication Set --- OpenClaw agent (independent, free-endpoints only)}
\date{2026-07-05}
\begin{document}
\maketitle

\section{Paper summary}

Freedman, Larsen, and Wang (2000) prove that certain \emph{topological modular functors} (TMFs) arising from Witten--Chern--Simons theory are \emph{universal for quantum computation}. Concretely, they show that any polynomial-time quantum circuit computation can be efficiently approximated by an intertwining action of a braid group representation on the state space of the SU(2) Chern--Simons TMF at a fifth root of unity (the ``CS5'' model, level $k=3$).

The technical heart of the paper is a \emph{density theorem}: the image of the Jones representation of the braid group $B_n$ (restricted to appropriate irreducible sectors) is dense in the special unitary group $\mathrm{SU}(V)$ modulo center, at $q = e^{2\pi i/5}$. Combined with the Solovay--Kitaev construction, this implies that any single-qubit or two-qubit target unitary can be approximated within accuracy $\varepsilon$ by a braid word of length $\mathrm{polylog}(1/\varepsilon)$.

The paper's abstract claim (\S1) is:
\begin{quote}
    \emph{A quantum circuit computation can be efficiently approximated by an intertwining action of a braid on the [modular] functor's state space $\ldots$ [the] computational model based on Chern--Simons theory at a fifth root of unity is $\ldots$ polynomially equivalent to the quantum circuit model.}
\end{quote}

Because this is a \emph{mathematical universality proof}, there is no headline empirical number to reproduce in the QC-100/QC-200 sense. The wave brief accordingly designates \textbf{SPOT-CHECK} as the appropriate verdict for this class of paper, provided the concrete substrate is verified numerically at small instance size. We do so by building the Fibonacci-anyon subcategory of SU(2)$_3$ (which is where the paper's density argument runs) and demonstrating:
\begin{enumerate}
    \item Its $F$ and $R$ symbols satisfy the pentagon and hexagon axioms (making it a well-defined unitary modular tensor category).
    \item The braid generators $\sigma_1, \sigma_2$ on the 2-dimensional 3-anyon computational subspace are unitary.
    \item Short braid words in $\{\sigma_1^{\pm 1}, \sigma_2^{\pm 1}\}$ come within operator-norm distance $\sim 3 \times 10^{-2}$ of the Hadamard and $T$ gates at braid length $\le 14$, numerically confirming the density claim on a small instance.
\end{enumerate}

\section{Claims table}
\begin{center}
\begin{tabular}{@{}llll@{}}
\toprule
\textbf{ID} & \textbf{Claim} & \textbf{Testable?} & \textbf{Tested?}\\
\midrule
C1 & Fibonacci $F$-symbol satisfies pentagon axiom & Yes (algebraic) & Yes (residual $\sim 10^{-16}$) \\
C2 & Fibonacci $(F,R)$ satisfies hexagon axiom (equivalently Yang--Baxter for $\sigma_1,\sigma_2$) & Yes (algebraic) & Yes (residual $\sim 10^{-16}$) \\
C3 & Braid representation $\rho: B_n \to U(V)$ is unitary & Yes (algebraic) & Yes (residual $\sim 10^{-16}$) \\
C4 & Braid words are dense in $\mathrm{SU}(V)$ (universality) & Testable at small $\varepsilon$ & Yes: dist to Hadamard $= 2.92\times 10^{-2}$ at length 13; dist to $T$ $= 3.19\times 10^{-2}$ at length 14 \\
C5 & Polynomial equivalence between CS5 and QCM & Not testable at small instance & No (asymptotic complexity claim) \\
C6 & Density of Jones representation in irreducible sectors (Theorem 1.1) & Very hard (algebraic geometry / Lie group density) & No (mathematical proof; we only verify concrete small-instance corollary) \\
\bottomrule
\end{tabular}
\end{center}

\section{Method (numbered)}

\subsection{Environment}
\begin{itemize}
    \item Host: CherryRd, macOS Darwin 25.3.0 (x86\_64), Python 3.14.6, NumPy 2.4.3, SciPy 1.18.0.
    \item No third-party TQFT/anyon library required; everything is implemented from scratch in pure NumPy so the reproduction is auditable and free of hidden dependencies.
    \item Endpoints: no LLM inference was used to produce the numerical results; the search is deterministic.
\end{itemize}

\subsection{Fibonacci anyon data}
Labels $\{0,1\}$ with $0\equiv\mathbf{1}$ (vacuum), $1\equiv\tau$. Fusion rule $\tau\otimes\tau=\mathbf{1}\oplus\tau$. The only non-trivial $F$-symbol is
\[
F^{\tau\tau\tau}_\tau \;=\;
\begin{pmatrix} 1/\varphi & 1/\sqrt{\varphi} \\ 1/\sqrt{\varphi} & -1/\varphi \end{pmatrix}, \qquad \varphi = \tfrac{1+\sqrt5}{2}.
\]
$R$-symbols:
\[
R^{\tau\tau}_{\mathbf{1}} = e^{-4\pi i / 5}, \qquad R^{\tau\tau}_{\tau} = e^{+3\pi i / 5}.
\]

\subsection{Computational qubit encoding}
On 3 $\tau$-anyons with total charge $\tau$, the fusion space is 2-dimensional with basis
$\ket{0}\equiv((\tau\tau)_\mathbf{1}\tau)_\tau,\;\;\ket{1}\equiv((\tau\tau)_\tau\tau)_\tau$.
The braid generators are:
\[
\rho(\sigma_1)=\operatorname{diag}(R^{\tau\tau}_\mathbf{1},R^{\tau\tau}_\tau),\qquad
\rho(\sigma_2)=F\cdot \operatorname{diag}(R^{\tau\tau}_\mathbf{1},R^{\tau\tau}_\tau)\cdot F.
\]

\subsection{Axiom checks}
Pentagon reduces (for Fibonacci) to $F^2 = I$ on the all-$\tau$ vertex; verified in code via \verb|check_pentagon_matrix|. Hexagon is verified equivalently via the Yang--Baxter braid relation $\sigma_1\sigma_2\sigma_1 = \sigma_2\sigma_1\sigma_2$ in \verb|check_hexagon|; unitarity of $\sigma_1,\sigma_2$ is verified in \verb|check_unitarity|.

\subsection{Braid-word density search}
We enumerate braid words of length $\le L$ in the four elementary generators $\{\sigma_1,\sigma_1^{-1},\sigma_2,\sigma_2^{-1}\}$, skip immediate cancellations, and canonicalise by a coarse fingerprint mod global phase to deduplicate the group element. For each word we compute the phase-invariant operator-norm distance
\[
d(U,V) = \min_\theta \|U - e^{i\theta}V\|_{\mathrm{op}}
\]
to a target gate $V$. We report the best distance found at each length.

\subsection{Reproducibility commands}
\begin{lstlisting}[language=bash]
cd work
python3 fibonacci_anyons.py --max-len 15 --out fibonacci_results.json
\end{lstlisting}
Wall clock at max-len 15: $\sim$ 3 minutes on one CPU core.

\section{Results vs.\ paper}

\begin{center}
\begin{tabular}{@{}lll@{}}
\toprule
\textbf{Quantity} & \textbf{Paper (implied / classical)} & \textbf{This work} \\
\midrule
Pentagon residual (all-$\tau$) & 0 (exact) & $1.11\times 10^{-16}$ (float64) \\
Hexagon / Yang--Baxter residual & 0 (exact) & $1.76\times 10^{-16}$ (float64) \\
$\sigma_1$ unitarity residual & 0 (exact) & $1.19\times 10^{-17}$ (float64) \\
$\sigma_2$ unitarity residual & 0 (exact) & $2.22\times 10^{-16}$ (float64) \\
$\rho(\sigma_1)$ eigenvalues & $e^{-4\pi i/5}, e^{+3\pi i/5}$ & $-0.809 - 0.588i,\; -0.309 + 0.951i$ (agree) \\
Density: $\min_L d(\rho(w), H)$ & density $\Rightarrow \varepsilon \to 0$ as $L\to\infty$ & $2.92\times 10^{-2}$ at $L=13$ \\
Density: $\min_L d(\rho(w), T)$ & density $\Rightarrow \varepsilon \to 0$ as $L\to\infty$ & $3.19\times 10^{-2}$ at $L=14$ \\
\bottomrule
\end{tabular}
\end{center}

\paragraph{Best braid word for Hadamard (length 13):}
\begin{lstlisting}
s1' s2 s1' s2 s1' s2 s2 s2 s1' s2 s1' s2 s1'
\end{lstlisting}
gives $d = 2.915526 \times 10^{-2}$.

\paragraph{Best braid word for $T$-gate (length 14):}
\begin{lstlisting}
s2 s1' s1' s2 s1' s2 s2 s1' s2 s1' s1' s2 s1' s1'
\end{lstlisting}
gives $d = 3.185377 \times 10^{-2}$.

\section{What worked / what did not}

\paragraph{What worked.}
\begin{itemize}
    \item Both pentagon and hexagon axioms are numerically saturated to machine precision.
    \item $\sigma_1,\sigma_2$ are unitary to machine precision.
    \item The BFS density search produces braid words that come within $\sim 3\%$ operator-norm distance of the Hadamard and $T$ gates at length $\le 14$, empirically confirming the density claim at a small scale.
    \item Global-phase-invariant hashing of the fingerprint successfully deduplicates the search frontier: without it the search space is $4^L$; with it we recover a polynomially bounded set of \emph{distinct} SU(2)-mod-center elements at each depth.
\end{itemize}

\paragraph{What did not work / limitations.}
\begin{itemize}
    \item We did not attempt to reproduce the full Freedman--Larsen--Wang density theorem (Theorem 1.1) or its Lie-group-theoretic proof. Only its small-instance numerical corollary is verified.
    \item The BFS search does not implement Solovay--Kitaev refinement; achieving $\varepsilon < 10^{-3}$ would require SK or a larger BFS depth than fits in this timebox.
    \item We only test the 1-qubit computational subspace ($n=3$ anyons, $\dim V = 2$). The paper's argument extends to $n=6$ anyons for the 2-qubit case (5-dimensional space with a 4-dimensional computational subspace); we did not build that here.
    \item Marker and Nougat CLIs are not installed on the host; the extraction files are honest \texttt{pdftotext} fallbacks (see \texttt{extraction/*} headers).
\end{itemize}

\section{Verdict}
\textbf{SPOT-CHECK} --- consistent with the QC wave brief's designation for mathematical universality proofs. The concrete Fibonacci-anyon substrate underlying the paper's density claim is fully verified numerically at small instance size (pentagon/hexagon/unitarity to machine precision, braid-word approximation of Hadamard and $T$ to $\sim 3\%$). We do not claim to have re-proven the density theorem itself; but the numerical evidence at $L\le 14$ is fully consistent with density and provides an operational witness of universality on a real anyon system.

\section*{Open Questions}

\textbf{Q1.} What is the empirical scaling of $\varepsilon_{\min}(L)$ for Hadamard/$T$ under Fibonacci braiding, and does it match the Solovay--Kitaev prediction $\varepsilon \sim 2^{-c L^{1/\alpha}}$ with $\alpha = \log_2 5$? Our BFS gives only two data points; a systematic sweep across $L = 4, 6, 8, \ldots, 20$ with efficient meet-in-the-middle deduplication would test the scaling law directly.

\textbf{Q2.} How does the numerical density rate compare on Fibonacci vs.\ Ising (SU(2)$_2$) vs.\ metaplectic (SU(2)$_4$) anyons at matched braid length? Ising is famously \emph{not} universal (Clifford only) --- can that be seen numerically as a plateau in $\varepsilon_{\min}(L)$?

\textbf{Q3.} The paper's 2-qubit gate construction uses $B_6$ on 6 anyons in the $V^0_{6\text{pts}} \oplus V^2_{6\text{pts}}$ space. What is the shortest braid word that approximates $\mathrm{CNOT}$ within $\varepsilon = 10^{-2}$? Is the ``leakage'' out of the computational subspace controllable, or does it accumulate to dominate the error budget?

\textbf{Q4.} Our density search deduplicates by a coarse-grained fingerprint of the SU(2) matrix rounded to $10^{-4}$. This gives an implicit lattice cover but may miss ``near-degenerate'' distinct elements that get quantized to the same bucket. Would a hyperbolic-geometry-aware nearest-neighbour index (using the natural Fubini--Study metric on $\mathrm{SU}(2)/\mathrm{U}(1) \simeq S^2$) find shorter braid approximations of the same targets?

\textbf{Q5.} Freedman--Larsen--Wang work at $r = 5$ (Fibonacci); Freedman--Kitaev--Wang (\emph{Density of the modular functor}, cs.CC/0006012) work more generally at $r \ge 5, r \ne 6$. Is the density-search efficiency (ratio of accessible distinct group elements to $4^L$) monotone in $r$, or does it peak at some intermediate value? A comparative BFS at $r = 5, 7, 9$ would give quantitative insight into what makes Fibonacci ``the best'' anyon for gate compilation.

\end{document}
